\cxtitle{An oblivious subspace injection need not give relative-error sketch-and-solve}

\begin{cxcredits}
\posedby{Amsel et al. \citeyearpar{AmselEtAl2026}}
\foundby{OpenAI Codex}{2026-08}
\formalizedby{OpenAI Codex}
\auditedby{OpenAI Codex}
\contributedby{Suvrit Sra}
\end{cxcredits}

\begin{cxcontext}
For a random matrix $\Omega\in\mathbb R^{n\times k}$ and a fixed
$r$-dimensional subspace $\mathcal V\subseteq\mathbb R^n$, the source calls
$\Omega$ a subspace injection with injectivity $\alpha$ on $\mathcal V$ when
$\mathbb E\|\Omega^Tx\|_2^2=\|x\|_2^2$ for every $x\in\mathbb R^n$ and
$\|\Omega^Tx\|_2^2\geq\alpha\|x\|_2^2$ for every $x\in\mathcal V$.
It is oblivious with failure probability $\delta$ if this lower inequality
holds on every fixed $r$-dimensional subspace with probability at least
$1-\delta$.  When no failure probability is displayed, the source fixes a
small constant, specifically $\delta=0.01$ as its example.  The discussion
immediately following Problems~5.1--5.2 asks whether constant OSI failure
probability suffices and then separately asks whether reducing it to
$\operatorname{poly}(\varepsilon)$ rescues the result.  Thus the probabilistic
claim tested here is the natural one in that context: the relative-error event
should hold with at least the OSI success probability.  Townsend and Wang
subsequently gave related counterexamples and resolved the question negatively
\cite{TownsendWang2026}.
\end{cxcontext}

\begin{cxsource}{osi-sketch-and-solve}
Problem~5.1 of the Simons workshop report of Amsel et al.
\cite{AmselEtAl2026}, scribed for Section~5 by Raphael A. Meyer.  Townsend and
Wang \cite{TownsendWang2026} had already resolved the general question; this
witness is independent and claims no priority.
\end{cxsource}

\begin{cxstatement}{osi-sketch-and-solve}
Let $A\in\mathbb{R}^{n\times d}$ be a full-rank matrix,
$B\in\mathbb{R}^{n\times p}$ be a matrix, and let
$\Omega\in\mathbb{R}^{n\times k}$ be an oblivious subspace injection with
dimension $d$ and injectivity $1-\varepsilon$. Let
$\widetilde{X}=(\Omega^{\top}A)^{+}(\Omega^{\top}B)\in\mathbb{R}^{d\times p}$
be the sketch-and-solve approximation to the ordinary least-squares solution
$X=A^{+}B$. (Here, ${}^{+}$ denotes the Moore--Penrose pseudoinverse.) Prove
or disprove: Does it necessarily hold that
\[
\|A\widetilde{X}-B\|_{\rm F}\leq(1+\mathcal{O}(\varepsilon))
\min_{X\in\mathbb{R}^{d\times p}}\|AX-B\|_{\rm F}?
\]
\end{cxstatement}

\begin{cxsummary}{osi-sketch-and-solve}
A one-dimensional oblivious subspace injection with injectivity $1$ can make
sketch-and-solve incur the fixed factor $\sqrt2$ with probability $1/50$.
\end{cxsummary}

\begin{cxcertificate}{osi-sketch-and-solve}
Three exact rationally weighted $2\times2$ sketch atoms, with exact isotropy,
an exact $99/100$ injection guarantee, and exact squared residual ratios.
\end{cxcertificate}

\begin{cxrefutation}
Problem~5.1 asks whether one-sided injectivity near $1$ yields the usual
relative-error guarantee for sketch-and-solve \cite{AmselEtAl2026}.  The
following three-atom distribution gives a negative answer under the failure
probability convention fixed in the source.

\begin{theorem}[\statusfalse: OSI relative error for sketch-and-solve]
\label{thm:osi-sketch-and-solve}
There is a $(1,1,1/100)$ oblivious subspace injection for which the claimed
relative-error inequality fails with probability $1/50$.  Consequently no
$1+O(\varepsilon)$ factor can hold with the OSI success probability, even
though the injectivity is exactly $1$.
\end{theorem}
\begin{proof}
Put $\rho=1/100$ and define
\[
 I=\begin{pmatrix}1&0\\0&1\end{pmatrix},\qquad
 B_+=\begin{pmatrix}1&0\\1&0\end{pmatrix},\qquad
 B_-=\begin{pmatrix}1&0\\-1&0\end{pmatrix}.
\]
Let $\Omega$ equal $I,B_+,B_-$ with respective probabilities
$1-2\rho,\rho,\rho$.  Its isotropy is exact:
\[
 \mathbb E[\Omega\Omega^T]
 =(1-2\rho)I+\rho
 \begin{pmatrix}1&1\\1&1\end{pmatrix}
 +\rho\begin{pmatrix}1&-1\\-1&1\end{pmatrix}=I.
\]
For $v=(x,y)^T$, the three possible values of
$\|\Omega^Tv\|_2^2$ are
\[
 x^2+y^2,\qquad (x+y)^2,\qquad (x-y)^2.
\]
The identity
$(x+y)^2+(x-y)^2=2(x^2+y^2)$ shows that at least one of the
last two values is at least $x^2+y^2$.  Hence, on every fixed line, the
injectivity-$1$ inequality can fail on at most one of the two atoms $B_+$ and
$B_-$, each of probability $\rho$.  Thus $\Omega$ is a
$(1,1,\rho)$-OSI, and in particular a $(1,1-\varepsilon,\rho)$-OSI for every
$\varepsilon\in(0,1)$.

Now take the single-response problem
\[
 A=\begin{pmatrix}1\\0\end{pmatrix},\qquad
 b=\begin{pmatrix}0\\1\end{pmatrix}.
\]
The unsketched optimum is $x_\star=0$ and its residual norm is $1$.  On the
identity atom, $\widetilde x=0$.  On the other atoms,
\[
 B_+^TA=\begin{pmatrix}1\\0\end{pmatrix},\quad
 B_+^Tb=\begin{pmatrix}1\\0\end{pmatrix},\qquad
 B_-^TA=\begin{pmatrix}1\\0\end{pmatrix},\quad
 B_-^Tb=\begin{pmatrix}-1\\0\end{pmatrix}.
\]
Therefore $\widetilde x=1$ on $B_+$ and $\widetilde x=-1$ on $B_-$, and in
both cases
\[
 \|A\widetilde x-b\|_2^2=\widetilde x^2+1=2.
\]
The fixed loss $\sqrt2$ occurs with total probability $2\rho=1/50$, whereas
the OSI failure probability is only $\rho=1/100$.

Finally, this is genuinely an asymptotic refutation of $1+O(\varepsilon)$.
If a universal constant $C>0$ supplied the hidden constant, choose
$\varepsilon>0$ so small that $C\varepsilon\leq1/4$.  The same sketch still
has injectivity at least $1-\varepsilon$, while on an event of probability
$1/50$ its ratio is $\sqrt2>5/4\geq1+C\varepsilon$.  No rounding or numerical
approximation is used.
\end{proof}
\end{cxrefutation}
